\section{Weekly Summary}

\begin{tcolorbox}[colback=yellow!10,colframe=black!50,title=AI-Generated Content]
\noindent The summary below was written by Claude (Anthropic) from that week's regression outputs and series descriptions. It has not been reviewed or fact-checked by a human, and should be read as an automated, informal recap -- not analysis.
\end{tcolorbox}

Welcome back to the weekly parade of statistical nonsense, where a robot pulls two economic time series out of a hat and demands they explain each other. As always, the correct amount of meaning to extract from any of this is zero, but we will spend several paragraphs extracting it anyway because that is the entire bit.

Monday kicked things off by pairing Coinbase Bitcoin with the San Francisco Tech Pulse index, and the raw regression came in absolutely radiant -- an R-squared near 0.70 and a p-value with more zeros in front of it than I care to count. "Bitcoin explains tech vibes!" screamed nobody responsible. Then we detrended it, and the R-squared collapsed to a limp 0.14. It survived, technically, in the sense that a soufflé survives being dropped, but that gorgeous raw fit was mostly two things drifting upward in the same era holding hands. Classic trend mirage, exactly the kind of thing this project exists to point at and laugh.

Then the week decided to humble us. Tuesday's Philadelphia capital-expenditures diffusion index versus the M2M user cost index produced a p-value of 0.89 -- the statistical equivalent of a shrug so total it echoes. Wednesday's Italian central bank forex interventions versus real-estate loan delinquencies flashed a "significant" raw result (p around 0.04) that promptly evaporated to 0.35 once detrended, another mirage caught red-handed. Thursday paired a recession indicator with the current account balance and got an R-squared of 0.005, which is roughly the correlation between my horoscope and my day. And Saturday's exports-to-newly-industrialized-countries versus state-and-local government's share of GDP gave us the week's spiciest twist: significant in the raw AND significant detrended, except the slope flips sign between the two. Truly nothing says "genuine relationship" like a coefficient that cannot decide which direction it points.

The two standouts actually worth a raised eyebrow: Friday, where U.S. CPI for nondurables squared off against French dwelling construction starts, held up after detrending (p around 0.04) -- weak, but it didn't just vanish. And Sunday's forestry-and-fishing corporate dividends versus a discontinued CPI variant posted a glorious raw R-squared of 0.85, and even after ripping out the shared trend it clung to a respectable 0.36 and stayed significant. Do net dividends from America's fishing sector secretly track medical-care-excluding inflation in "Size Class D"? Almost certainly not, and I beg you not to trade on it. These are randomly paired series with no theory, no causal story, and no statistical rigor beyond "the line goes through the dots-ish." Spurious correlation is the house specialty here, not the exception -- so enjoy the pretty numbers, and then immediately forget them. See you next week for more.
